% =====================================================================
%  FICHE 8 — THÈME 2 — NIVEAU DIFFICILE — CORRIGÉ
% =====================================================================
\documentclass[11pt,a4paper]{article}
\usepackage[utf8]{inputenc}
\usepackage[T1]{fontenc}
\usepackage{lmodern}
\usepackage[french]{babel}
\usepackage{amsmath,amssymb,mathtools}
\usepackage{icomma}
\usepackage{siunitx}
\sisetup{output-decimal-marker={,},per-mode=power,group-separator={\,},group-minimum-digits=5}
\usepackage[version=4]{mhchem}
\usepackage{xcolor}
\usepackage{tikz}
\usepackage[most]{tcolorbox}
\usepackage{enumitem}
\usepackage{array}
\usepackage{fancyhdr}
\usepackage[hidelinks]{hyperref}
\usetikzlibrary{arrows.meta,positioning,calc}
\usepackage[a4paper,margin=2.1cm,headheight=26pt,headsep=10pt]{geometry}
\usepackage{titlesec}

\definecolor{primary}{RGB}{150,98,162}
\definecolor{primarydark}{RGB}{92,52,110}
\definecolor{excol}{RGB}{0,135,90}

\tcbset{boxbase/.style={enhanced,breakable,boxrule=0.9pt,arc=2.5pt,left=3mm,right=3mm,top=2.3mm,bottom=2.3mm,fonttitle=\bfseries,coltitle=white,toptitle=0.6mm,bottomtitle=0.6mm}}
\newtcolorbox[auto counter]{correction}[1][]{boxbase,colback=excol!4,colframe=excol,title={\textsc{Corrigé}~\thetcbcounter\ \ifx\relax#1\relax\relax\else\textmd{--- \itshape #1}\fi}}

\newcommand{\dH}{\Delta H}
\newcommand{\dHr}{\Delta H^{\circ}_{\text{r}}}
\newcommand{\dHf}{\Delta H^{\circ}_{\text{f}}}
\newcommand{\rep}[1]{\textcolor{primarydark}{\textbf{#1}}}

\pagestyle{fancy}
\fancyhf{}
\fancyhead[L]{\small\textcolor{primarydark}{\textbf{Fiche 8} — Thème 2 : Enthalpie de formation et loi de Hess}}
\fancyhead[R]{\small\textcolor{primarydark}{\textsc{Difficile — corrigé}}}
\fancyfoot[C]{\small\thepage}
\renewcommand{\headrulewidth}{0.8pt}
\renewcommand{\headrule}{\hbox to\headwidth{\color{primary}\leaders\hrule height \headrulewidth\hfill}}
\setlength{\parindent}{0pt}\setlength{\parskip}{3pt}

\begin{document}

\begin{center}
\begin{tikzpicture}
\node[rectangle,rounded corners=7pt,fill=excol,text=white,minimum width=16cm,
      minimum height=1.1cm,align=center,inner sep=6pt]
  {{\large\bfseries Thème 2 — Enthalpie de formation et loi de Hess}\\[1pt]{\small Niveau \textsc{Difficile} — corrigé détaillé}};
\end{tikzpicture}
\end{center}

\begin{correction}[enthalpie de formation du méthane par combinaison]
\begin{enumerate}[label=\alph*),leftmargin=2em,topsep=2pt,itemsep=3pt]
  \item Formation à partir des corps simples : \;$\ce{C}(s) + 2\,\ce{H2}(g) \rightarrow \ce{CH4}(g)$.
  \item Combinaison : $(1) + 2\times(2) - (3)$. En effet, en inversant $(3)$ et en ajoutant $(1)$ et
        $2\times(2)$ :
  \begin{align*}
    &\ce{C} + \ce{O2} \rightarrow \ce{CO2} & (\dH_1)\\
    &2\,\ce{H2} + \ce{O2} \rightarrow 2\,\ce{H2O} & (2\dH_2)\\
    &\ce{CO2} + 2\,\ce{H2O} \rightarrow \ce{CH4} + 2\,\ce{O2} & (-\dH_3)
  \end{align*}
  On simplifie $\ce{CO2}$, $2\,\ce{H2O}$ et $2\,\ce{O2}$ : il reste $\ce{C} + 2\,\ce{H2} \rightarrow \ce{CH4}$. \checkmark
  \item $\dHf(\ce{CH4}) = \dH_1 + 2\,\dH_2 - \dH_3 = -394 + 2(-286) - (-890) = -394 - 572 + 890 = \rep{\SI{-76}{\kilo\joule\per\mol}}$.
        Valeur négative : la formation du méthane est \rep{exothermique}.
\end{enumerate}
\end{correction}

\begin{correction}[cycle énergétique de la formation du dioxyde d'azote]
\begin{enumerate}[label=\alph*),leftmargin=2em,topsep=2pt,itemsep=2pt]
  \item Depuis les éléments, deux chemins mènent à $2\,\ce{NO2}$ : directement ($+66$), ou via
        $2\,\ce{NO}+\ce{O2}$ ($+180$) puis $\dH$. Égalité des chemins :
        \[ (+180) + \dH = (+66). \]
  \item $\dH = 66 - 180 = \rep{\SI{-114}{\kilo\joule}}$ pour $2\,\ce{NO}+\ce{O2}\rightarrow 2\,\ce{NO2}$.
  \item $\dH<0$ : cette étape est \rep{exothermique}. La réaction globale
        $\ce{N2}+2\,\ce{O2}\rightarrow2\,\ce{NO2}$ a $\dH=+66>0$ : elle est \rep{endothermique}
        (former $\ce{NO}$ à partir de $\ce{N2}$ coûte beaucoup d'énergie).
  \item Par mole de $\ce{NO2}$ : $\dfrac{-114}{2} = \rep{\SI{-57}{\kilo\joule\per\mol}}$.
\end{enumerate}
\end{correction}

\begin{correction}[réduction de l'oxyde de fer par le carbone]
\begin{enumerate}[label=\alph*),leftmargin=2em,topsep=2pt,itemsep=3pt]
  \item Il faut \textbf{inverser $(1)$ et la multiplier par $2$} (pour obtenir $2\,\ce{Fe2O3}$ à gauche)
        et \textbf{multiplier $(2)$ par $3$} :
  \begin{align*}
    &2\,\ce{Fe2O3}(s) \rightarrow 4\,\ce{Fe}(s) + 3\,\ce{O2}(g) & (-2\dH_1)\\
    &3\,\ce{C}(s) + 3\,\ce{O2}(g) \rightarrow 3\,\ce{CO2}(g) & (3\dH_2)
  \end{align*}
  En additionnant, les $3\,\ce{O2}$ se simplifient : $\;2\,\ce{Fe2O3} + 3\,\ce{C} \rightarrow 4\,\ce{Fe} + 3\,\ce{CO2}$. \checkmark
  \item $\dH = -2\,\dH_1 + 3\,\dH_2 = -2(-824) + 3(-394) = 1648 - 1182 = \rep{\SI{+466}{\kilo\joule}}$.
  \item $\dH>0$ : la réduction est \rep{endothermique}. C'est pourquoi le haut fourneau doit être
        \emph{fortement chauffé} : il faut apporter de l'énergie (par la combustion du coke) pour
        arracher l'oxygène au minerai de fer.
\end{enumerate}
\end{correction}

\end{document}
